\section{Hệ thống đề xuất}
\subsection{Mô tả tổng quan hệ thống}
Hệ thống được giới thiệu trong luận văn này có tên là Food Tour Planner, là
một ứng dụng web nhằm hỗ trợ người dùng dễ dàng tìm kiếm thông tin về các quán ăn, xây dựng và quản
lý các hành trình ăn uống (food tour), đồng thời xác định lộ trình di chuyển hợp lý và tiết kiệm thời gian giữa nhiều
địa điểm ẩm thực trong cùng một hành trình. Hệ thống được thiết kế nhằm cung cấp một nền
tảng tích hợp, giúp người dùng lập kế hoạch ăn uống một cách thuận tiện, hiệu quả và có khả
năng chia sẻ với cộng đồng.

Các chức năng chính của hệ thống bao gồm:

\begin{itemize}[label=\textbullet]
    \item \textbf{Quản lý người dùng:} Đăng ký/Đăng nhập và quản lý thông tin tài khoản người dùng.
    \item \textbf{Quản lý và tra cứu thông tin quán ăn:} Tìm kiếm và xem thông tin chi tiết
    về các quán ăn, bao gồm vị trí và các thông tin liên quan phục vụ nhu cầu tham khảo.
    \item \textbf{Đóng góp thông tin nhà hàng:} Người dùng có thể gửi yêu cầu đăng tải thông tin nhà hàng mới. Các quán này
    được công khai sau khi quản trị viên phê duyệt.
    \item \textbf{Tạo và quản lý food tour:} Tạo các hành trình food tour cá nhân
    bằng cách thêm một hoặc nhiều quán ăn vào cùng một lộ trình.
    \item \textbf{Sắp xếp lộ trình di chuyển:} Xác định thứ tự ghé thăm các quán ăn trong food tour
    nhằm tối ưu hóa quãng đường và thời gian di chuyển.
    \item \textbf{Gợi ý food tour:} Đề xuất các food tour dựa trên mức độ tương đồng giữa các quán ăn trong
    lịch sử sử dụng của người dùng và các food tour đã tồn tại.
    \item \textbf{Chia sẻ food tour công khai:} Công khai food tour để người dùng khác
    có thể tìm kiếm, tham khảo và xem lại.
\end{itemize}

Hệ thống được thiết kế để đáp ứng nhu cầu của các nhóm người dùng chính sau:

\begin{itemize}[label=\textbullet]

    \item \textbf{Người dùng}: 
    Đây là nhóm người dùng chính. Họ dùng ứng dụng để tìm kiếm quán ăn, 
    tạo và quản lý các food tour cá nhân, sắp xếp lộ trình đi lại, đánh giá 
    quán ăn và chia sẻ food tour với cộng đồng. Ngoài ra, người dùng có thể 
    tìm kiếm và học hỏi từ các food tour đã được chia sẻ công khai để xây dựng 
    hành trình ăn uống phù hợp với sở thích riêng.

    \item \textbf{Người dùng quản lý nhà hàng (Restaurant Owner)}:
    Là nhóm người dùng có nhu cầu đưa thông tin nhà hàng của mình lên hệ thống.
    Nhóm người dùng này vẫn sử dụng tài khoản người dùng thông thường nhưng được
    cấp thêm quyền quản lý nhà hàng. Người dùng quản lý nhà hàng có thể gửi yêu
    cầu đăng tải thông tin nhà hàng mới, cung cấp các thông tin liên quan như tên
    nhà hàng, địa chỉ, loại hình ẩm thực, hình ảnh và mô tả. Các thông tin nhà
    hàng chỉ được hiển thị công khai sau khi được quản trị viên kiểm tra và
    phê duyệt.

    \item \textbf{Quản trị viên hệ thống (System Administrator)}:
    Chịu trách nhiệm kiểm duyệt các yêu cầu đăng tải nhà hàng, quản lý và vận hành hệ thống, bao gồm quản lý tài khoản
    người dùng, quản lý dữ liệu nhà hàng, giám sát nội dung food tour công khai và đảm bảo hệ thống hoạt động
    ổn định, an toàn.

\end{itemize}

Với kiến trúc hướng dịch vụ và tập trung vào trải nghiệm người dùng, 
hệ thống được xây dựng để có thể dễ dàng được mở rộng cùng như bảo trì và 
linh hoạt trong việc phát triển thêm các tính năng mới trong tương lai.

\subsection{Stakeholder}
Các stakeholder là những cá nhân hoặc nhóm có liên quan trực tiếp hoặc gián tiếp đến hệ
thống, có lợi ích bị ảnh hưởng hoặc tham gia vào quá trình phát triển, vận hành và sử dụng hệ
thống. Dựa trên mức độ liên quan, các stakeholder của hệ thống được phân thành các nhóm
chính và phụ như sau.

Các stakeholder chính (Primary stakeholders):
\begin{itemize}[label=\textbullet]
    \item \textbf{Người dùng cuối (End Users)}: Nhóm người dùng sử dụng trực tiếp hệ thống để tìm kiếm thông tin quán ăn, tạo và quản lý các food tour cá nhân, sắp xếp lộ trình di chuyển và chia sẻ hoặc tham khảo các food tour công khai. Đây là nhóm stakeholder chính, quyết định mức độ thành công và hiệu quả sử dụng của hệ thống.
    \item \textbf{Người dùng quản lý nhà hàng (Restaurant Owner)}:
    Nhóm người dùng muốn đưa thông tin nhà hàng lên hệ thống.
    Nhóm này sử dụng các chức năng gửi/yêu cầu đăng tải nhà hàng,
    cập nhật thông tin nhà hàng và theo dõi trạng thái xử lý.
    Các nhà hàng do nhóm người dùng này gửi chỉ được công khai trên hệ thống
    sau khi được quản trị viên kiểm tra và phê duyệt.
    \item \textbf{Quản trị viên hệ thống (System Administrator)}:
Là người chịu trách nhiệm quản lý và vận hành hệ thống, bao gồm quản lý tài khoản người dùng,
quản lý dữ liệu quán ăn, giám sát nội dung food tour được chia sẻ công khai và đảm bảo hệ thống hoạt động ổn định, an toàn.
\end{itemize}

Các stakeholder phụ (Secondary stakeholders): 

\begin{itemize}[label=\textbullet]
    \item \textbf{Nhóm phát triển hệ thống (Development Team)}: Là nhóm thiết kế và triển khai hệ thống, chịu trách nhiệm phân tích yêu cầu, thiết kế kiến trúc, xây dựng cơ sở dữ liệu và phát triển các chức năng của ứng dụng.
\end{itemize}



\subsection{Yêu cầu chức năng (Functional Requirements)}
\textbf{FR1: Quản lý tài khoản người dùng}
\begin{itemize}[label=\textbullet]
    \item Cung cấp các chức năng đăng ký tài khoản, đăng nhập, đăng xuất và quản lý mật khẩu
nhằm đảm bảo người dùng có thể truy cập và sử dụng hệ thống.
    \item Người dùng phải có khả năng xem và cập nhật các thông tin hồ sơ cá nhân cơ bản của mình.
\end{itemize}

\textbf{FR2: Quản lý và kiểm duyệt hệ thốn(Quản trị viên)}
\begin{itemize}[label=\textbullet]

    \item Cung cấp chức năng cho quản trị viên quản lý
    dữ liệu nền tảng, bao gồm quản lý tài khoản người dùng,
    quản lý các nhà hàng đã được phê duyệt và giám sát nội dung
    được chia sẻ trên hệ thống.

    \item Quản trị viên phải có khả năng xem danh sách các yêu cầu
    đăng tải nhà hàng từ người dùng có quyền quản lý nhà hàng.

    \item Quản trị viên có thể kiểm tra thông tin nhà hàng,
    thực hiện phê duyệt hoặc từ chối yêu cầu đăng tải.

    \item Các nhà hàng chỉ được hiển thị công khai trên hệ thống
    sau khi được quản trị viên phê duyệt.

\end{itemize}

\textbf{FR3: Đăng ký và quản lý thông tin nhà hàng}

\begin{itemize}[label=\textbullet]

    \item Người dùng có quyền quản lý nhà hàng có thể gửi yêu cầu
    đăng tải nhà hàng mới lên hệ thống.

    \item Yêu cầu đăng tải nhà hàng phải bao gồm các thông tin cơ bản:
    tên nhà hàng, địa chỉ, loại hình ẩm thực, hình ảnh,
    mô tả và các thông tin liên quan.

    \item Hệ thống phải lưu trạng thái xử lý của yêu cầu đăng tải,
    bao gồm: Pending, Approved, Rejected.

    \item Người dùng có quyền quản lý nhà hàng có thể theo dõi
    trạng thái các yêu cầu đã gửi.

\end{itemize}

\textbf{FR4: Tìm kiếm và tra cứu thông tin quán ăn/nhà hàng} 
\begin{itemize}[label=\textbullet]
    \item Người dùng có thể tìm kiếm quán ăn theo nhiều tiêu chí khác nhau như tên
quán ăn, ẩm thực, địa điểm, giá cả, thời gian mở cửa, đánh giá và khoảng cách.
    \item Kết quả tìm kiếm được hiển thị dưới dạng danh sách kèm theo 
    thông tin tổng quan và cho phép người dùng xem chi tiết từng quán ăn.
    \item Sắp xếp kết quả tìm kiếm theo các tiêu chí 
    phổ biến như khoảng cách, mức giá và mức độ đánh giá.
\end{itemize}

\textbf{FR5: Quản lý quán ăn yêu thích} 
\begin{itemize}[label=\textbullet]
    \item Thêm hoặc bỏ quán ăn ra khỏi danh sách yêu thích cá
nhân.
    \item Xem và quản lý danh sách các quán ăn đã được đánh dấu yêu thích.
\end{itemize}

\textbf{FR6: Tạo và quản lý Food Tour cá nhân} 
\begin{itemize}[label=\textbullet]
    \item Người dùng tạo mới food tour với các thông tin cơ bản như tiêu đề và
mô tả.
    \item Thêm hoặc xóa các quán ăn vào food tour và sắp xếp thứ tự các điểm ăn uống trong hành trình.
    \item Cập nhật và chỉnh sửa thông tin food tour 
    của người dùng trong quá trình sử dụng.
\end{itemize}

\textbf{FR7: Tối ưu lộ trình Food Tour} 
\begin{itemize}[label=\textbullet]
    \item Từ danh sách các điểm ăn uống trong food tour, hệ thống hỗ trợ sắp xếp lộ trình di chuyển
hợp lý theo yêu cầu của người dùng nhằm giảm tổng quãng đường di chuyển.
    \item Hiển thị tổng khoảng cách di chuyển của food tour sau khi thực hiện tối ưu.
    \item Người dùng có thể thực hiện lại việc tối ưu lộ trình khi danh sách hoặc thứ tự các
điểm dừng thay đổi.
\end{itemize}

\textbf{FR8: Chia sẻ Food Tour (Public/Private)} 
\begin{itemize}[label=\textbullet]
    \item Thiết lập trạng thái hiển thị của food tour ở chế độ riêng
tư (private) hoặc công khai (public).
    \item Các food tour được chia sẻ ở chế độ công khai phải có thể được người dùng khác tìm kiếm và
xem nội dung chi tiết.
\end{itemize}

\textbf{FR9: Tìm kiếm và xem Food Tour public} 
\begin{itemize}[label=\textbullet]
    \item Tìm kiếm các food tour công khai dựa trên từ khóa như tên
    tour hoặc mô tả.
    \item Xem chi tiết food tour công khai, bao gồm danh sách quán ăn, 
    thứ tự ghé thăm và tổng khoảng cách di chuyển.
\end{itemize}

\textbf{FR10: Lưu và yêu thích Food Tour public} 
\begin{itemize}[label=\textbullet]
    \item Lưu food tour công khai của người dùng khác để tham
khảo và sử dụng sau này.
    \item Đánh giá cụ thể mức độ yêu thích đối với food tour công khai, và hệ
thống phải hiển thị số lượt yêu thích trung bình của mỗi food tour.
\end{itemize}

\textbf{FR11: Đánh giá quán ăn (Review)} 
\begin{itemize}[label=\textbullet]
    \item Hệ thống cho phép người dùng tạo đánh giá cho quán ăn dưới dạng nội dung nhận xét và
điểm đánh giá.
    \item Các đánh giá của người dùng phải được hiển thị trong trang thông tin chi tiết của quán ăn để
phục vụ cho việc tham khảo của cộng đồng.
\end{itemize}

\textbf{FR12: Gợi ý quán ăn theo khẩu vị cá nhân} 
\begin{itemize}[label=\textbullet]
    \item Lưu trữ thông tin khẩu vị cá nhân của người dùng như một phần của hồ sơ
người dùng.
    \item Dựa trên mức độ tương đồng giữa khẩu vị người dùng và đặc điểm ẩm thực của quán ăn, hệ
thống đề xuất các quán ăn phù hợp nhằm cá nhân hóa trải nghiệm tìm kiếm.
    \item Lưu trữ thêm thông tin khu vực ưu tiên và lịch sử food tour của người dùng
     để tăng độ chính xác của gợi ý.
    \item Hiển thị rõ ràng lý do gợi ý cho từng quán ăn (ví dụ: “Bạn thích món cay • Đã 
    xuất hiện trong 2 food tour trước của bạn”).
    \item Hỗ trợ gợi ý quán ăn ngay trong quá trình tạo food tour để người dùng dễ dàng bổ sung 
    điểm dừng.
    \item Kết quả gợi ý có thể được sắp xếp theo độ phù hợp, khoảng cách hoặc đánh giá.
\end{itemize}

\textbf{FR13: Hiển thị bản đồ và lộ trình FoodTour} 
\begin{itemize}[label=\textbullet]
    \item Hệ thống cung cấp chức năng hiển thị bản đồ tương tác cho mỗi food tour.
    \item Trên bản đồ phải thể hiện vị trí các quán ăn, thứ tự ghé thăm và lộ trình di chuyển giữa các
điểm dừng.
    \item Bản đồ được cập nhật tương ứng khi người dùng thay đổi danh sách hoặc thứ tự các điểm
dừng trong food tour.
\end{itemize}


\subsection{Yêu cầu phi chức năng (Non-functional Requirements)}

\textbf{NFR1: Hiệu năng (Performance)} 
\begin{itemize}[label=\textbullet]
    \item Hệ thống phải đảm bảo thời gian phản hồi nhanh dưới 2s trong các thao tác chính như tìm kiếm quán
ăn, tìm kiếm food tour và xem thông tin chi tiết, nhằm không làm gián đoạn trải nghiệm người
dùng.
    \item Chức năng hiển thị bản đồ và lộ trình di chuyển phải hoạt động mượt mà, đảm bảo khả năng
tương tác liên tục khi người dùng thay đổi lộ trình.
    \item Hiệu năng của hệ thống phải được duy trì ổn định khi số lượng người dùng dưới 50 người.
\end{itemize}

\textbf{NFR2: Độ tin cậy và tính sẵn sàng (Reliability \& Availability)} 
\begin{itemize}[label=\textbullet]
    \item Hệ thống phải đảm bảo độ ổn định cao và có khả năng hoạt động liên tục trong điều kiện sử
dụng thông thường.
    \item Hệ thống có cơ chế xử lý lỗi và ghi nhận log để thuận tiện cho việc khắc phục
\end{itemize}

\textbf{NFR3: Bảo mật (Security)} 
\begin{itemize}[label=\textbullet]
    \item Thông tin người dùng, bao gồm thông tin cá nhân, khẩu vị và lịch sử sử dụng, phải được lưu
trữ và bảo vệ an toàn.
    \item Thông tin đăng nhập của người dùng phải được mã hóa và không lưu trữ dưới dạng rõ.
    \item Hệ thống phải có cơ chế kiểm soát truy cập, đảm bảo người dùng chỉ có thể thao tác trên các
tài nguyên phù hợp với quyền hạn của mình.
    \item Các food tour ở chế độ riêng tư (private) chỉ được phép truy cập bởi chủ sở hữu.
\end{itemize}

\textbf{NFR4: Tính khả dụng (Usability)} 
\begin{itemize}[label=\textbullet]
    \item Giao diện người dùng phải được thiết kế trực quan, dễ hiểu và dễ sử dụng, phù hợp với cả
những người dùng không có nhiều kiến thức kỹ thuật.
    \item Các chức năng chính như tìm kiếm quán ăn, tạo food tour và chia sẻ food tour phải được
trình bày rõ ràng, giúp người dùng dễ dàng thao tác và làm quen với hệ thống.
\end{itemize}

\textbf{NFR5: Khả năng mở rộng và bảo trì (Scalability \& Maintainability)} 
\begin{itemize}[label=\textbullet]
    \item Hệ thống phải được thiết kế theo hướng module hóa, giúp dễ dàng mở rộng hoặc bổ sung
chức năng trong tương lai.
    \item Cấu trúc hệ thống phải hỗ trợ việc bảo trì, chỉnh sửa và nâng cấp mà không ảnh hưởng lớn
đến các chức năng hiện có.
\end{itemize}